% ============================================================================
% Main-text section "The Audit on a Second Platform" (new Section 5, placed
% after sec04_validation_submission). Every numeric claim is bound to a
% \valFedAud* / \valExp* / \valArmor* macro defined in values.tex (verified
% cell-by-cell against outputs/comprasnet/...; see
% docs/jleo_rr_revision/NEW_NUMBERS_MAP_COMPRASNET.md). Voice, macro style, and
% table conventions mirror sec04_validation_submission.tex and
% sec05_cobidder_type_submission.tex.
% ============================================================================

\section{The Audit on a Second Platform}
\label{sec:comparative}

The validation in Section~\ref{sec:validation} draws a boundary on
one platform: under a reproducible, non-circular label, the
award-layer score concentrates adjudication-anchored loser-side
exposure in the raw data, but almost nothing survives once
procurement opportunity is held fixed, and the prospective ranking
collapses outside the incumbent pool. The defensible product is not
the screen but the audit---the decomposition that tells an agency
how far a cheap award-layer statistic can be trusted before costly
bid-layer work begins. A decomposition is only worth carrying,
however, if it travels. This section runs the same audit, unchanged,
on a second procurement system, and asks what the protocol returns
where the institutions differ. The exercise tests portability of the
\emph{method} and of the loser-side construct; it is not an
independent cartel universe, and we read it as neither confirmation
nor refutation of the legal status of any firm.

\subsection{A Different Procurement System}
\label{sec:comparative_setting}

The federal platform differs from BEC--SP on the dimensions that
matter for the audit. ComprasNet is a single national system: where
BEC is one state-level platform inside the federal legal frame,
ComprasNet aggregates procuring units (UASGs) across the whole
federal administration. Its modality mix is also narrower. The
federal panel is effectively pure Preg\~ao---about $15\%$ regular
electronic auction and $85\%$ price-registration (SRP, \emph{sistema
de registro de pre\c{c}os})---and the sealed-bid Convite modality
that supplies the institutional contrast in the main text is
extinct federally. The audit therefore runs on a panel with no
modality margin to stratify on, which sharpens rather than weakens
the test: the loser-side construct must earn its keep without the
modality variation that organized the BEC reading. The observable
window is $2013$--$2019$, the federal-overlap window with the BEC
panel.

The observability gap is the second contrast, and it is on-thesis.
The public federal data expose participation and the winner flag,
but no bid microdata---there is no federal analogue to the BEC
\texttt{LANCES} bid ladder. A screen that reads only the award layer
is exactly the screen one \emph{can} build here; the bid-layer
forensic step that the main text positions downstream of the screen
cannot be constructed from the public federal release at all. The
federal panel is thus the natural stress test for a cheap award-layer
audit: it is the environment where bid-microdata screens are
unavailable, so the question of how far the award layer alone can be
trusted is not academic but binding.

\subsection{Data, Label, and Linkage}
\label{sec:comparative_data}

The federal panel comprises $\valFedPanelRows$ participation rows
across $\valFedFirms$ distinct CNPJ-14 firms,
built from the Comptroller-General (CGU) open bulk dumps and
restricted to the Preg\~ao modalities (regular and SRP) under the
canonical schema used throughout the main text. Of these,
$\valFedAlwaysLosers$ are always-losers---firms with a win rate of zero over
the observed window. Applying the \emph{same} IQR rule used for
BEC---$\textit{median} + 1.5\cdot\textit{IQR}$ on the tenders-count
distribution of always-losers, not the Tukey upper-quartile
rule---re-derives a federal frequent-loser threshold of $\valFedFLthreshold$
(against $14$ on BEC), and the canonical federal frequent-loser
count above that threshold is $\valFedFLcount$ firms. The threshold is
re-estimated from the federal data, not transported: the construct
is the rule, not the cutoff value, exactly as the rules of
engagement in the main text require.

The legal anchors are the \emph{same} CADE cases that anchor the BEC
validation---$\valFedAudCases$ numbered adjudicated procurement-cartel
processes, with one unnumbered summary case excluded for lack of
firm-level identification. This is the honest statement of what the
federal exercise is: a second platform with \emph{partially
overlapping} legal anchors, not a fresh cartel population. The
federal evidence therefore tests whether the audit protocol and the
loser-side construct \emph{port} across procurement systems; it does
not assemble an independent ground truth, and we do not read it as
one. Direct defendants are matched to federal participant records by
the canonical establishment-anchored CNPJ rule. Expanding the
anchored tender-items yields $\valFedAudCobiddersAll$ federal cobidders under the
broad rule and $\valFedAudCobiddersBroadAL$ under the broad always-loser restriction
that mirrors the main-text validation target; the frequent-loser
flag is never used to build the label, so the federal validation is
non-circular in the score by the same construction as the BEC
validation.

\subsection{Side-by-Side Audit Verdict}
\label{sec:comparative_verdict}

Table~\ref{tab:comparative_audit} runs each stage of the
Section~\ref{sec:validation} battery on both platforms and reports
the comparable quantity in each column. The point of the table is
not the absolute levels but the \emph{shape} of the audit: whether
each successive discipline---opportunity adjustment, within-stratum
restriction, power-bounding, matched permutation, label-frozen
timing, case concentration, negative controls---moves the federal
verdict the way it moves the BEC verdict.

\begin{table}[htbp]
\centering
\caption{The Audit on Two Platforms: BEC--SP versus ComprasNet}
\label{tab:comparative_audit}
\scriptsize
\begin{adjustbox}{max width=\textwidth,center}
\begin{threeparttable}
\begin{tabular}{p{5.4cm}cc}
\toprule
Audit stage & BEC--SP & ComprasNet \\
\midrule
Raw award-layer ROC-AUC & $\valExpRawScoreAUC$ & $\valFedAudRawAUC$ \\
\quad Exposure-only ROC-AUC (no score) & $\valExpOnlyAUC$ & $\valFedAudExpOnlyAUC$ \\
Label-blind opportunity expectation (AUC) & $\valArmorExpLB$ & $\valFedAudLabelBlindAUC$ \\
Within-stratum residual (AUC, MEDIUM cells) & $\valExpWithinAUC$ & $\valFedAudWithinAUC$ \\
Nested DeLong increment over exposure-only ($p$) & $+\valExpIncrement$ ($p=\valExpDeLongP$) & $\valFedAudNestedIncrement$ ($p=\valFedAudNestedP$) \\
Power-bounded residual (det.\ prob.\ @ true AUC $0.55$ / $0.60$)\tnote{a} & $\valArmorPowerFive$\,/\,$\valArmorPowerTen$ & $\valFedAudPowerFive$\,/\,$\valFedAudPowerSixty$ \\
Matched-permutation $p$ (within-strata shuffle) & $\valExpPermCp$ & $\valFedAudPermCp$ \\
Label-frozen prospective timing (AUC)\tnote{b} & $\valArmorFrozenProspAUC$ & $\valFedAudFrozenProspAUC$ \\
Top-case concentration (share of positives) & $\valTopCasePosShare$ & $\valFedAudTopCaseShare$ \\
Negative-control verdict & generic geometry & \valFedAudNegControlVerdict \\
\bottomrule
\end{tabular}
\begin{tablenotes}
\scriptsize
\item[a] The federal positive set is roughly $30\%$ of BEC's, so the
matched-permutation null is informative against within-stratum
residuals $\geq 0.60$ (detection probability $\valFedAudPowerSixty$) but
underpowered against a residual of $0.55$ (detection probability
$\valFedAudPowerFive$); BEC reaches $\valArmorPowerFive$ already at $0.55$. The federal
within-stratum positive control $O_i$ recovers $\geq 0.99$
($\valFedAudOiControlMedium$), so the
design provably detects within-stratum signal when present---the
federal null is a power bound on a genuine null, not a design
artifact.
\item[b] The federal prospective cell freezes score and always-loser
pool and labels new defendant contact in $2017$--$2019$
(positives $N^{+}=\valFedAudFrozenProspN$); the retrospective companion
($\valFedAudFrozenRetroAUC$, $N^{+}=\valFedAudFrozenRetroN$) scores in-window contact under the
same frozen pool. The two are different estimands on different
positive sets and are mutually consistent, not contradictory; only
the prospective number is the referee's clean out-of-time test.
\item \textit{Notes:} Each row is the comparable quantity from the
Section~\ref{sec:validation} battery, re-run on the federal panel
under the same definitions; the federal construction is detailed in
\ref{app:federal_audit_submission}. Cobidder labels are
adjudication-anchored exposure, not cartel membership, on both
platforms; neither column ranks defendants. The federal panel has no
sealed-bid modality, so no modality stratification is reported.
The decisive row is the exposure-only line directly beneath raw
AUC: where the opportunity axis alone matches (or, federally,
edges out) the raw score, the raw number is measuring procurement
opportunity, not conduct. The raw-AUC and timing rows measure
different estimands on different samples---raw award-layer
discrimination over the full panel versus label-frozen prospective
ranking on rankable incumbents---and must not be collided. PR-AUC is
the lead metric for the rare-target rows; ROC-AUC is shown for
comparability, and PR-AUC levels must be read against each platform's
base rate---the federal positive rate is roughly an order of
magnitude lower than BEC's, so federal PR-AUC is bounded below BEC's
by the rarer target, not by weaker discrimination.
\end{tablenotes}
\end{threeparttable}
\end{adjustbox}
\end{table}

\subsection{Interpretation}
\label{sec:comparative_interpretation}

On the federal panel the audit returns the same verdict it returns
on BEC, and on the row that matters it returns it more cleanly. On
both platforms the raw validation over-credits the screen: the raw
award-layer ROC-AUC ($\valFedAudRawAUC$ federally) looks like discrimination
until one asks what an analyst would have predicted from procurement
opportunity alone. Federally that question answers itself in a
single line---the exposure-only ranking, which never sees the score,
reaches $\valFedAudExpOnlyAUC$, matching and if anything edging out the raw
score itself ($\valFedAudRawAUC$). This is the cleanest demonstration yet
that the raw number measures who shows up, not who colludes: where on
BEC the exposure axis only approaches the raw score, federally it
already equals it, so the screen's apparent reach is opportunity
arithmetic with no conduct premium left to claim. The rest of the
battery confirms the reading from below. Once procurement
opportunity is held fixed the within-stratum residual collapses to
near chance ($\valFedAudWithinAUC$), the nested model adds nothing detectable
over the exposure-only benchmark ($\valFedAudNestedIncrement$ DeLong increment,
$p=\valFedAudNestedP$), and the matched-strata permutation does not reject
($p=\valFedAudPermCp$)---the same collapse, the same nulls, the same shape
the BEC audit produces. We state the federal power bound as a bound,
not as a result it cannot bear: with a federal positive set about a
third the size of BEC's, the matched-permutation null is informative
against a within-stratum residual of $0.60$ but underpowered at
$0.55$ ($\valFedAudPowerSixty$ versus $\valFedAudPowerFive$ detection probability), so the
federal deflation does not rest on the permutation alone. It rests on
four power-robust pillars that no sample-size argument can dissolve:
the exposure-only ranking already matches the raw score, the
within-stratum point estimate sits \emph{below} $0.5$, the matched
permutation does not reject, and the negative controls return the
same generic geometry; and the within-stratum positive control
recovers $\geq 0.99$ ($\valFedAudOiControlMedium$), confirming the design would have
caught a within-stratum signal had one been present. The deflation is
therefore not a federal artifact and not a BEC artifact; it is what a
disciplined audit
returns whenever a cheap award-layer statistic is held to an
opportunity-adjusted standard. That is the argument, not a
confession: what ports across the two systems is the audit---the
decomposition that catches the over-crediting wherever it is
run---not a deployable loser-side ranking. An agency on either
platform that runs this decomposition learns, before opening any bid
record, that loser-side concentration flags opportunity and should
not be carried as a stand-alone prioritization device; and on the
federal platform that lesson binds harder, because the bid-layer step
that might follow the screen cannot be constructed from public data
at all.

\subsection{Observability as an On-Thesis Finding}
\label{sec:comparative_observability}

The two platforms differ in what they let an analyst see, and the
differences are not nuisance---they are the finding. They differ
first in scale of the target: the federal positive base rate is
roughly an order of magnitude lower than BEC's,
so any cross-platform PR-AUC comparison must be read against
the base rate that bounds it, not as a difference in screen quality:
a rarer target mechanically depresses precision-recall area, and the
federal PR-AUC sits below BEC's for that reason before any question
of discrimination arises. Item groups are
not observed federally in the BEC sense: federal item codes are
embedded in buyer-specific catalogs rather than a shared
classification, so the item-group margin that the BEC audit
stratifies on (the COARSE cell) has no clean federal counterpart,
and the federal opportunity cells are built on the buyer$\times$year
margins instead (\ref{app:federal_audit_submission}). The bid tier
is absent entirely: no federal bid ladder means no bid-dispersion
benchmark, no bid-layer forensic comparator, and no way to run the
downstream costly-proof step on public data. Both gaps point the
same way. The award-layer audit is precisely the audit that remains
feasible where the richer information is missing, which is the
common case for an enforcement agency staring at a national platform
with participation records and little else. That the loser-side
construct can be re-derived, the label re-anchored, and the full
deflation battery re-run on a system with no bid microdata is itself
the portability result: the cheap screen runs where the expensive
screens cannot be built, and the audit tells the agency how far to
trust it before the expensive screens---if they ever become
available---are worth commissioning.
